\documentclass[12pt]{article}

\usepackage{amsmath}
\usepackage{german}
\usepackage{graphicx}
\DeclareMathOperator\sgn{sgn}

%\graphicspath{{/home/ritz/Teaching/fcn/fig/}}

%----------------- layout ----------------------------------------

\oddsidemargin=0in
\textwidth=6in
\topmargin=-0.7in
\textheight=9.5in
\baselineskip=16pt
\font\mittel=cmr17
\font\roman=cmr12
\font\klein=cmr9
\font\mgross=cmr10
\newfont{\CAP}{cmcsc10 scaled 900}
\pagestyle{empty}
\parindent=0mm

%----------------- macros ----------------------------------------
\newcommand{\conc}[1]{[{\rm C}]_{\text{#1}}}
\newcommand{\units}[1]{\text{#1}}
\newcommand{\K}{\text{K}}
\newcommand{\Na}{\text{Na}}
\newcommand{\Cl}{\text{Cl}}
\newcommand{\Leak}{\text{leak}}
\newcommand{\water}{\text{H}_2\text{O}}
\newcommand{\tderiv}[1]{\frac{\mathrm{d}#1}{\mathrm{d}t}}
%-----------------------------------------------------------------


\begin{document}

\parbox{2cm}{
\includegraphics[width=1.8cm]{../bccnlogo.pdf}
}
\parbox{11cm}{
\begin{center}
\large HUMBOLDT-UNIVERSIT"AT \hskip 0.1 cm ZU \hskip 0.1 cm BERLIN
\vskip 0.1 true cm
\mgross BERNSTEIN CENTER FOR\\ COMPUTATIONAL NEUROSCIENCE
\end{center}
}
\parbox{2cm}
{
\hfill
\includegraphics[width=1.8cm]{../hublogo.pdf}
}

\vskip 0.6 true cm

\leftline{\CAP Humboldt-Universit"at zu Berlin 
\hfill Phone: 030/2093-9110}
\leftline{\CAP Philippstr. 13 House 6
             \hfill Fax: 030/2093-6771}
\leftline{\CAP \hfill webpage: http://www.bccn-berlin.de/}


\vskip 0.8 true cm
\centerline{\bf Models of Neural Systems, WS 2010/11}
\centerline{\bf Project N/A: Short-term synaptic plasticity} 
\centerline{Project presentation and report submission: N/A}

\vskip 0.8 true cm

\section{Background}

\section{}
\begin{enumerate}
    \item Tsodyks-Markram model of short-term synaptic depression        

        \begin{enumerate}
            \item implement the model of short-term synaptic
                depression by Tsodyks and Markram \cite{},
            \item using short spike trains with pre-define inter-spike
                intervals show the effects of depression and recovery,
            \item demonstrate how depression and recovery depend on
                parameters of the model: $\tau_{\mathrm{rec}}$,
                $U_{\mathrm{SE}}$ and $\tau_{\mathrm{inact}}$.
        \end{enumerate}

    \item Limiting frequency

        \begin{enumerate}
            \item feed the dynamic synapse with an input modelled with
                a Poisson spike train with a pre-defined firing rate
                and no refractoriness
            \item plot the steady-state amplitude of excitatory post-synaptic
                current  (EPSC) as a function of the input firing rate
            \item demonstrate that over a certain input frequency
                (limitting frequency) the output firing rate saturates
                and transmits no more information about the input firing
                rate
        \end{enumerate}

    \item Regulating neural code

        \begin{enumerate}
            \item simulate the responses of the neuron to inputs with
            constant firing rates and fast-transient
        \item investigate in which parameter regimes the synapse
            transmits information about the mean firing rate and  the
            transients
        \end{enumerate}

    \item Gain modulation

        This part of the project may be solved in one of simulator of
        spiking neural networks (PyNN or Brian).

The goal of this part is to show that synaptic depression may
underline multiplication of neuronal signals \cite{}.

\begin{enumerate}
\item simulate a leaky integrate-and-fire neuron with inhibitory and
    excitatory synpases, where only excitatory synapses are depressing

\item plot an output firing rate as a function of the input firing
    rate at different levels of inhibition. show that inhibition
    modulates the input-output gain.

\item repeat the above experiment without synaptic depression. How
    does the input-output relation change?
\end{enumerate}


\end{enumerate}



\vfill
\begin{tabular}{lll}
Urs Bergmann & Phone: 2093-8924 & Email:
urs.bergmann(at)biologie.hu-berlin.de \\
Matthias Guggenmos & & Email: matthias.guggenmos(at)bccn-berlin.de \\
Richard Kempter \hfill & Phone: 2093-8925 \hfill & Email:
r.kempter(at)biologie.hu-berlin.de \\
\end{tabular}

\end{document}
